%% toc: P3 30 30.1,30.2,30.3,30.4
\section{Proofs of the substructure, extension and chain preservation theorems}
\label{sec:p03-w03-p03-026-preservation-proofs}

This section proves, in both directions, the three syntactic
characterizations stated in \cref{sec:p03-preservation}: the class form of
the \L o\'s--Tarski theorem (\cref{thm:p03-los-tarski}, closure under
substructures against universal axioms), its sentence form (preservation
under extensions against existential sentences), and the
Chang--\L o\'s--Suszko theorem (\cref{thm:p03-chang-los-suszko}, closure under
unions of arbitrary chains of substructures against \(\forall\exists\)
axioms).  The positive-formula theorem for surjective homomorphisms is
proved separately in \cref{sec:p03-n-a-011-02}.

\emph{Inputs.}  The arguments use four results already proved in the
handbook: compactness of first-order logic
(\cref{thm:fol-compactness}; it is applied below to signatures
\(\Lang(M)\) that are uncountable whenever \(M\) is, so the converse
directions use the Axiom of Choice assumed there), the
diagram tests for embeddings and elementary embeddings
(\cref{thm:diagram-tests}), invariance of satisfaction under
isomorphism (\cref{lem:iso-invariance}) and the elementary chain
theorem (\cref{thm:elementary-chain}).  Nothing else is imported.

\emph{Conclusions.}  Each characterization is proved over the class of all
nonempty \(\Lang\)-structures, for an arbitrary signature \(\Lang\) with
relation, operation and constant symbols and with logical equality.  The
class forms are proved relative to an ambient set of axioms, which is the
generality of \L o\'s's 1955 Theorems~5 and~6
\cite[pp.~47--48]{Los1955}; the absolute statements of
\cref{sec:p03-preservation} are the case of an empty ambient set.

\subsection{Prefix classes, consequence sets and chains}
\label{sec:p03-w03-p03-026-conventions}

\begin{hbconvention}{def:p03-w03-p03-026-prefix-classes}{Universal, existential and \(\forall\exists\) sentences}
A formula is \emph{quantifier-free} if it is built from atomic formulas,
equality atoms included, by \(\neg,\land,\lor,\to\).  For a quantifier-free
\(\theta\) and finite (possibly empty) variable tuples \(\bar x,\bar y\):
\begin{itemize}
 \item a \emph{universal} sentence has the form \(\forall\bar x\,\theta(\bar x)\);
 \item an \emph{existential} sentence has the form \(\exists\bar x\,\theta(\bar x)\);
 \item a \emph{\(\forall\exists\) sentence} has the form
       \(\forall\bar x\,\exists\bar y\,\theta(\bar x,\bar y)\).
\end{itemize}
An empty block is permitted, so a quantifier-free sentence is universal,
existential and \(\forall\exists\), and every universal or existential
sentence is a \(\forall\exists\) sentence.  The sentence \(\forall x\,(x=x)\)
is used as the empty conjunction and \(\exists x\,\neg(x=x)\) as the empty
disjunction. On nonempty domains they are respectively equivalent to
\(\exists x\,(x=x)\) and \(\forall x\,\neg(x=x)\), so truth and falsity
each have representatives in all three prefix classes.  Structures have nonempty
domains.  For a set \(T\) of \(\Lang\)-sentences put
\[
 T_{\forall}=\{\psi:\psi\text{ universal and }T\entails\psi\},\qquad
 T_{\forall\exists}=\{\psi:\psi\text{ a }\forall\exists\text{ sentence and }T\entails\psi\}.
\]
Thus \(T\entails T_{\forall\exists}\entails T_{\forall}\), and
\(\operatorname{Mod}(T)\subseteq\operatorname{Mod}(T_{\forall\exists})\subseteq\operatorname{Mod}(T_{\forall})\).
\end{hbconvention}

\begin{hbdefinition}{def:p03-w03-p03-026-chain}{Chains of substructures and their unions}
Let \(\delta\) be a nonzero ordinal.  A \emph{chain} is a family
\(\langle\mathcal M_i:i<\delta\rangle\) of \(\Lang\)-structures with
\(\mathcal M_i\subseteq\mathcal M_j\) whenever \(i<j<\delta\).  Its
\emph{union} \(\bigcup_{i<\delta}\mathcal M_i\) has domain
\(\bigcup_{i<\delta}\dom\mathcal M_i\); a relation symbol is interpreted by
the union of its interpretations, an operation symbol by the common value of
its interpretations at any stage containing the arguments, and a constant by
its common value.  An \emph{elementary chain}
(\cref{sec:maps-elementarity-05}) is the special case in which every
inclusion is elementary.  Closure of a class under unions of chains, as in
\cref{def:p03-preservation-relations}, refers to arbitrary chains of
substructures.
\end{hbdefinition}

For a structure \(\mathcal A\), \(\operatorname{Diag}(\mathcal A)\) and
\(\operatorname{ElDiag}(\mathcal A)\) are the atomic and elementary diagrams
of \cref{def:diagrams}, in the signature \(\Lang(A)\) of
\cref{def:parameter-expansion}.  Write
\(\operatorname{Diag}_{\forall}(\mathcal A)\) for the set of universal
\(\Lang(A)\)-sentences true in \(\mathcal A_A\).  It contains
\(\operatorname{Diag}(\mathcal A)\), since a literal is a universal sentence
with an empty block.

\begin{hblemma}{thm:p03-w03-p03-026-union}{proved=annex:p03-w03-p03-026-union}{The union of a chain is a common extension}
For every chain \(\langle\mathcal M_i:i<\delta\rangle\) the union
\(\mathcal M\) is a well-defined \(\Lang\)-structure, and
\(\mathcal M_i\subseteq\mathcal M\) for every \(i<\delta\).
\end{hblemma}

\begin{hblemma}{thm:p03-w03-p03-026-easy}{proved=annex:p03-w03-p03-026-easy}{The syntactic sides are preserved}
Let \(\mathcal M\subseteq\mathcal N\).
\begin{enumerate}
 \item A quantifier-free \(\theta(\bar x)\) satisfies
       \(\mathcal M\entails\theta(\bar a)\iff\mathcal N\entails\theta(\bar a)\)
       for every tuple \(\bar a\) from \(M\).
 \item If \(\mathcal N\entails\forall\bar x\,\theta(\bar x,\bar a)\) with
       \(\theta\) quantifier-free and \(\bar a\) from \(M\), then
       \(\mathcal M\entails\forall\bar x\,\theta(\bar x,\bar a)\).  Dually,
       existential sentences with parameters from \(M\) pass from
       \(\mathcal M\) to \(\mathcal N\).
 \item If every member of a chain satisfies a \(\forall\exists\) sentence,
       so does the union of the chain.
\end{enumerate}
\end{hblemma}

\begin{hblemma}{thm:p03-w03-p03-026-prenex}{proved=annex:p03-w03-p03-026-prenex}{Closure of the prefix classes}
The negation of a universal sentence is logically equivalent to an
existential sentence, and conversely.  A finite conjunction of universal
sentences is logically equivalent to a universal sentence, and a finite
disjunction of existential sentences to an existential sentence.  The same
holds for sentences with parameters, uniformly in the parameters.
\end{hblemma}

\subsection{The \L o\'s--Tarski theorem}
\label{sec:p03-w03-p03-026-los-tarski}

The converse directions rest on one embedding criterion.  It identifies the
substructures of models of \(T\) by the universal part of \(T\) alone.

\begin{hblemma}{thm:p03-w03-p03-026-embed}{proved=annex:p03-w03-p03-026-embed}{Substructures of models of a theory}
For a set \(T\) of \(\Lang\)-sentences and an \(\Lang\)-structure
\(\mathcal M\), the following are equivalent:
\begin{enumerate}
 \item \(\mathcal M\entails T_{\forall}\);
 \item \(\mathcal M\subseteq\mathcal N\) for some \(\mathcal N\entails T\).
\end{enumerate}
\end{hblemma}

\begin{hbtheorem}{thm:p03-w03-p03-026-los-tarski-relative}{proved=annex:p03-w03-p03-026-los-tarski-relative}{\L o\'s--Tarski theorem relative to an ambient theory}
Let \(S\subseteq T\) be sets of \(\Lang\)-sentences.  The following are
equivalent:
\begin{enumerate}
 \item every model of \(S\) that is a substructure of a model of \(T\) is a
       model of \(T\);
 \item there is a set \(U\) of universal sentences with
       \(\operatorname{Mod}(T)=\operatorname{Mod}(S\cup U)\).
\end{enumerate}
When these hold, \(U=T_{\forall}\) is such a set.
\end{hbtheorem}

\begin{hbcorollary}{thm:p03-w03-p03-026-los-tarski-class}{proved=annex:p03-w03-p03-026-los-tarski-class}{\L o\'s--Tarski theorem, class form}
For an elementary class \(\mathcal K=\operatorname{Mod}(T)\) of
\(\Lang\)-structures, the following are equivalent:
\begin{enumerate}
 \item \(\mathcal K\) is closed under substructures;
 \item \(\mathcal K\) is axiomatizable by universal \(\Lang\)-sentences.
\end{enumerate}
In that case \(\mathcal K=\operatorname{Mod}(T_{\forall})\).  This is the
class statement of \cref{thm:p03-los-tarski}.
\end{hbcorollary}

\begin{hbtheorem}{thm:p03-w03-p03-026-los-tarski-sentence}{proved=annex:p03-w03-p03-026-los-tarski-sentence}{\L o\'s--Tarski theorem, sentence form}
Let \(S\) be a set of \(\Lang\)-sentences and \(\varphi\) an
\(\Lang\)-sentence.
\begin{enumerate}
 \item The following are equivalent: (a) whenever
       \(\mathcal M\subseteq\mathcal N\) are models of \(S\) and
       \(\mathcal M\entails\varphi\), also \(\mathcal N\entails\varphi\);
       (b) there is an existential sentence \(\psi\) with
       \(S\entails\varphi\leftrightarrow\psi\).
 \item The following are equivalent: (a) whenever
       \(\mathcal M\subseteq\mathcal N\) are models of \(S\) and
       \(\mathcal N\entails\varphi\), also \(\mathcal M\entails\varphi\);
       (b) there is a universal sentence \(\chi\) with
       \(S\entails\varphi\leftrightarrow\chi\).
\end{enumerate}
For \(S=\varnothing\), item~1 is the sentence statement of
\cref{thm:p03-los-tarski}: a sentence is preserved under extensions between
all \(\Lang\)-structures iff it is logically equivalent to an existential
sentence.
\end{hbtheorem}

\L o\'s calls a sentence satisfying item~1(a) \emph{persistent} with respect
to the class \(\operatorname{Mod}(S)\); his Theorem~6 is the implication from
1(a) to 1(b), with the equivalence derivable from the axioms
\cite[p.~48]{Los1955}.  The ambient theory matters.  Every sentence is
persistent with respect to the class of algebraically closed fields, since
quantifier elimination for algebraically closed fields makes any inclusion
of algebraically closed fields elementary, as noted after the
Tarski--Vaught test in \cref{sec:maps-elementarity-03}.  Item~1 then yields, for every sentence, an
existential sentence equivalent to it modulo the axioms of algebraically
closed fields.  Over all structures in the ring language the same sentences
need not have existential equivalents: the sentence
\(\forall x\,\exists y\,(x=0\lor x\cdot y=1)\) holds in the field of
rationals and fails in its extension, the polynomial ring in one
indeterminate over the rationals, where the indeterminate has no inverse; by
\cref{thm:p03-w03-p03-026-easy} it is therefore equivalent to no existential
sentence over all structures.

\subsection{Unions of chains and the Chang--\L o\'s--Suszko theorem}
\label{sec:p03-w03-p03-026-chang-los-suszko}

The converse direction of the chain theorem builds an elementary chain and
an arbitrary chain with the same union.  Two lemmas supply the alternating
steps.

\begin{hblemma}{thm:p03-w03-p03-026-step-up}{proved=annex:p03-w03-p03-026-step-up}{A model of \(T\) above a model of \(T_{\forall\exists}\)}
Let \(\mathcal A\entails T_{\forall\exists}\).  There is
\(\mathcal B\entails T\) with \(\mathcal A\subseteq\mathcal B\) such that
every existential \(\Lang(A)\)-sentence true in \(\mathcal B_A\) is true in
\(\mathcal A_A\).
\end{hblemma}

\begin{hblemma}{thm:p03-w03-p03-026-step-back}{proved=annex:p03-w03-p03-026-step-back}{Returning to an elementary extension}
Let \(\mathcal A\subseteq\mathcal B\), and suppose every existential
\(\Lang(A)\)-sentence true in \(\mathcal B_A\) is true in \(\mathcal A_A\).
There is \(\mathcal C\) with \(\mathcal A\preccurlyeq\mathcal C\) and
\(\mathcal B\subseteq\mathcal C\).
\end{hblemma}

\begin{hbtheorem}{thm:p03-w03-p03-026-chang-los-suszko}{proved=annex:p03-w03-p03-026-chang-los-suszko}{Chang--\L o\'s--Suszko theorem}
For a set \(T\) of \(\Lang\)-sentences, the following are equivalent:
\begin{enumerate}
 \item \(\operatorname{Mod}(T)\) is closed under unions of chains of
       substructures, in the sense of \cref{def:p03-w03-p03-026-chain};
 \item \(\operatorname{Mod}(T)\) is closed under unions of chains indexed
       by \(\omega\);
 \item \(\operatorname{Mod}(T)=\operatorname{Mod}(T_{\forall\exists})\);
 \item \(\operatorname{Mod}(T)\) is axiomatizable by \(\forall\exists\)
       sentences \(\forall\bar x\,\exists\bar y\,\theta(\bar x,\bar y)\)
       with \(\theta\) quantifier-free, either block possibly empty.
\end{enumerate}
The equivalence of items~1 and~4 is \cref{thm:p03-chang-los-suszko}.
\end{hbtheorem}

\begin{hbtheorem}{thm:p03-w03-p03-026-elementary-chains-differ}{proved=annex:p03-w03-p03-026-elementary-chains-differ}{Elementary chains do not characterize}
Every elementary class is closed under unions of elementary chains.  The
class of linear orders with a greatest element is elementary, and it is not
closed under unions of chains of substructures; hence it is not
axiomatizable by \(\forall\exists\) sentences, and its closure under unions
of elementary chains does not imply the hypothesis of
\cref{thm:p03-w03-p03-026-chang-los-suszko}.
\end{hbtheorem}

\subsection{A worked example: groups in the multiplication signature}
\label{sec:p03-w03-p03-026-groups}

\L o\'s observes that a group described with only a binary multiplication
can have a substructure that is not a group, for instance the natural
numbers inside the additive integers, while the description with a unary
inverse operation admits open (quantifier-free) axioms
\cite[p.~47]{Los1955}.  The two preservation theorems sort out which syntactic
form each description can take.

\begin{hbcorollary}{thm:p03-w03-p03-026-groups}{proved=annex:p03-w03-p03-026-groups}{Groups with one binary operation}
In the signature \(\{\cdot\}\) with one binary operation symbol, let
\(\mathcal G\) be the class of structures \((G,\cdot)\) that are groups.
\begin{enumerate}
 \item \(\mathcal G=\operatorname{Mod}(\Gamma)\) for the two
       \(\forall\exists\) sentences
       \[
        \forall x\,y\,z\,\bigl((x\cdot y)\cdot z=x\cdot(y\cdot z)\bigr),
        \qquad
        \forall a\,b\,\exists x\,y\,(a\cdot x=b\land y\cdot a=b).
       \]
 \item \(\mathcal G\) is closed under unions of chains of substructures.
 \item \(\mathcal G\) is not closed under substructures, so it is not
       axiomatizable by universal sentences in this signature.
\end{enumerate}
\end{hbcorollary}

The example illustrates the procedure the theorems license.  To show that a
class has no universal axiomatization in a given signature, exhibit one
substructure of a member that lies outside the class; to show that it has a
\(\forall\exists\) axiomatization without finding one, prove the class
elementary and closed under unions of chains.  Item~2 can also be checked
directly from item~1 by \cref{thm:p03-w03-p03-026-easy}; the theorem adds the
converse, which turns a semantic closure argument into existence of axioms of
the required shape.

\begin{hbwarning}{def:p03-w03-p03-026-finite}{What the proofs use beyond finite structures}
The preserving directions (\cref{thm:p03-w03-p03-026-easy}) are proved
structure by structure and remain valid when every structure involved is
finite.  The converse directions are different.  The structures produced by
\cref{thm:p03-w03-p03-026-embed,thm:p03-w03-p03-026-step-up,thm:p03-w03-p03-026-step-back}
come from compactness, which fails when only finite structures are admitted:
the sentences ``there are at least \(n\) elements'', for \(n\ge1\), are
satisfiable together in no finite structure, although each finite subset is.
The chain condition also loses its content among finite structures: if a
chain of finite structures has a finite union, its domains form an increasing
sequence of subsets of a finite set, so some stage equals the union.  Every
class of finite structures is therefore closed under those unions of chains
that remain finite, and the hypothesis of
\cref{thm:p03-w03-p03-026-chang-los-suszko} carries no information there.
The statements of this section are therefore claims about the class of all
structures; \cref{sec:p03-preservation} records that the substructure form of
the \L o\'s--Tarski theorem fails when the structures are restricted to
finite ones.
\end{hbwarning}

\subsection{\L o\'s 1955: extension with primary and secondary conditions}
\label{sec:p03-w03-p03-026-los1955}

\L o\'s's paper \emph{On the extending of models (I)} \cite{Los1955} asks
when a model \(\mathfrak M\) of a set \(\mathcal A\) of axioms can be enlarged
to a model \(\mathfrak M^*\) in which \(\mathcal A\) still holds and every
sentence of a further set \(\mathcal C\) holds as well.  He calls the members
of \(\mathcal A\) \emph{primary conditions of extension} and those of
\(\mathcal C\) \emph{secondary conditions of extension}, and he seeks the
conditions on \(\mathfrak M\) that are necessary and sufficient for such an
extension to exist \cite[p.~38]{Los1955}.  His opening case is a commutative
ring that is to be extended to a field, where the familiar answer is the
absence of zero divisors; written as
\(x\cdot y=0\to(x=0\lor y=0)\), that condition is an open formula, with no
quantifiers \cite[pp.~38--39]{Los1955}.  The main result,
\cref{thm:p03-w03-p03-026-los55-thm1}, shows that this is the general
pattern: the necessary and sufficient conditions can always be taken to be a
set of open formulas \cite[p.~39]{Los1955}.  The paper was received on
1~February 1954 and was presented in part to the Toru\'n section of the
Polish Mathematical Society on 12~May 1953 and at the VIIIth Congress of
Polish Mathematicians in Warsaw in September 1953
\cite[pp.~38, 54]{Los1955}.

This subsection follows the paper section by section and keeps its
notation; inside this subsection \(\mathcal A,\mathcal B,\mathcal C,\mathcal X\)
denote sets of formulas, \(\mathfrak M\) a model and \(T\) an index set of
constants, as in the source.  Its \S\S\,8--9 are the relative forms of the
\L o\'s--Tarski theorem proved in
\cref{sec:p03-w03-p03-026-los-tarski} by the embedding criterion; here they
are derived by \L o\'s's own route, from the extension theorem.

\subsubsection{The three theories and their models (\S\S\,1--3)}

\begin{hbconvention}{def:p03-w03-p03-026-los55-theories}{The theories \(\mathcal E_1\subset\mathcal E_2\subset\mathcal E_3\)}
\L o\'s works with three elementary theories that differ only in their
primitive signs \cite[pp.~39--40]{Los1955}.  The formulas of
\(\mathcal E_1\) are built from relation signs \(r_1,r_2,\ldots\) of arities
\(\nu_1,\nu_2,\ldots\) and function signs \(f_1,f_2,\ldots\) of arities
\(\varkappa_1,\varkappa_2,\ldots\); \(\mathcal E_2\) adds relation signs
\(q_1,q_2,\ldots\) of arities \(\mu_1,\mu_2,\ldots\); \(\mathcal E_3\) adds
individual constants \(g_t\), \(t\in T\), for an arbitrary index set \(T\).
The identity sign \(=\) occurs in all three.  The connectives are \(\to\),
\(\lor\), \(\land\), \(\equiv\) and the postfix negation \('\); the
quantifiers are the general quantifier \(\prod_{x}\) and the existential
quantifier \(\sum_{x}\); the variables are \(x_1,x_2,\ldots\)
\cite[p.~40]{Los1955}.  Thus
\begin{equation*}
 \text{(1.1)}\qquad \mathcal E_1\subset\mathcal E_2\subset\mathcal E_3 .
\end{equation*}
Within \(\mathcal E_i\), \(\mathcal S_i\) is the set of sentences (no free
variables) and \(\mathcal O_i\) the set of open formulas (no bound
variables).  The two sets overlap: in \(\mathcal E_3\) a formula built from
constants alone, with no variables, is both \cite[p.~40]{Los1955}.

A \emph{model} for \(\mathcal E_i\) is a nonempty set \(A\) with relations
\(R_k\) of arity \(\nu_k\) for \(r_k\), functions \(F_k\colon
A^{\varkappa_k}\to A\) for \(f_k\), and, from \(\mathcal E_2\) on, relations
\(Q_k\) of arity \(\mu_k\) for \(q_k\), and in \(\mathcal E_3\) elements
\(G_t\in A\) for \(g_t\), not necessarily distinct \cite[p.~41]{Los1955}:
\(\mathfrak M=\langle A,R_1,R_2,\ldots,F_1,F_2,\ldots,Q_1,Q_2,\ldots,\{G_t\}_{t\in T}\rangle\).
A formula is \emph{valid} in \(\mathfrak M\) only when \(\mathfrak M\)
interprets each of its signs; a formula with free variables is valid iff its
universal closure is, so free variables behave as generally quantified ones;
\(=\) is always identity \cite[p.~42]{Los1955}.  \(\mathfrak M\) is a
\emph{model of} \(\mathcal X\) if every formula of \(\mathcal X\) is valid in
it.  In the handbook's terms, \(\mathcal E_i\) is the set of first-order
formulas over a signature \(\Lang_i\), a model for \(\mathcal E_i\) is an
\(\Lang_i\)-structure, and validity is truth of the universal closure.
\end{hbconvention}

\L o\'s uses the derivability relation of the functional calculus with
identity, writing \(\mathit{Cn}_i(\mathcal X)\) for the set of formulas of
\(\mathcal E_i\) derivable from \(\mathcal X\subset\mathcal E_i\); with the
rule of omitting and adding general quantifiers \cite[pp.~39--42]{Los1955},
a formula with free variables in \(\mathcal X\) acts as its universal
closure.  \(\mathcal X\) is \emph{consistent} in \(\mathcal E_i\) if
\(\mathit{Cn}_i(\mathcal X)\neq\mathcal E_i\) \cite[p.~41]{Los1955}.  The
paper treats the following properties as known; the handbook proves them
from soundness and strong completeness
(\cref{thm:fol-soundness,thm:fol-completeness}).

\begin{hblemma}{thm:p03-w03-p03-026-los55-cn}{proved=annex:p03-w03-p03-026-los55-cn}{The consequence operation and G\"odel's theorem}
Let \(\mathcal X\subset\mathcal E_i\).
\begin{enumerate}
 \item \(\mathit{Cn}_i(\mathcal X)\) is the set of \(\alpha\in\mathcal E_i\)
       valid in every model for \(\mathcal E_i\) of \(\mathcal X\).  Hence
       (2.1) a model of \(\mathcal X\) is a model of
       \(\mathit{Cn}_i(\mathcal X)\); (2.2) a set with a model is consistent;
       and G\"odel's theorem holds: every consistent set has a model, for
       any number of primitive signs, the Axiom of Choice being used when
       the signs are not denumerable \cite[p.~42]{Los1955}.
 \item (1.2) \(\mathcal X\subset\mathit{Cn}_i(\mathcal X)=\mathit{Cn}_i(\mathit{Cn}_i(\mathcal X))\subset\mathcal E_i\).
       (1.3) \(\mathit{Cn}_i(\mathcal X)\) is the union of the sets
       \(\mathit{Cn}_i(\mathcal Y)\) over finite \(\mathcal Y\subset\mathcal X\).
       (1.4) If \(\alpha\in\mathcal S_i\) and \(\beta\in\mathcal E_i\), then
       \(\beta\in\mathit{Cn}_i(\mathcal X\cup(\alpha))\) iff
       \(\alpha\to\beta\in\mathit{Cn}_i(\mathcal X)\).
       (1.5) Some \(\alpha\in\mathcal E_i\) has
       \(\mathit{Cn}_i((\alpha))=\mathcal E_i\).
       (1.6) An inconsistent \(\mathcal X\) has a finite inconsistent subset.
       (1.7) If \(\alpha\in\mathcal S_i\) and \(\mathcal X\cup(\alpha)\) is
       inconsistent, then \(\alpha'\in\mathit{Cn}_i(\mathcal X)\)
       \cite[pp.~40--41]{Los1955}.
 \item (1.8) If \(\alpha(g_{t_1},\ldots,g_{t_n})\in\mathcal S_3\) lists all
       constants of the sentence, \(\mathcal X\subset\mathcal E_2\), and
       \(\alpha(g_{t_1},\ldots,g_{t_n})\in\mathit{Cn}_3(\mathcal X)\), then
       \(\alpha(x_{s_1},\ldots,x_{s_n})\in\mathit{Cn}_2(\mathcal X)\) for
       any variables \(x_{s_1},\ldots,x_{s_n}\) not occurring in
       \(\alpha(g_{t_1},\ldots,g_{t_n})\), where
       \(\alpha(x_{s_1},\ldots,x_{s_n})\) replaces each \(g_{t_k}\) by
       \(x_{s_k}\) \cite[p.~41]{Los1955}.
\end{enumerate}
\end{hblemma}

\begin{hbdefinition}{def:p03-w03-p03-026-los55-extensions}{Submodels and the two kinds of extension}
Let \(\mathfrak M=\langle A,R_1,\ldots,F_1,\ldots\rangle\) and
\(\mathfrak M^*=\langle A^*,R^*_1,\ldots,F^*_1,\ldots\rangle\) be models for
\(\mathcal E_1\).  If \(A\subset A^*\), \(R^*_k|A=R_k\) and \(F^*_k|A=F_k\)
for every \(k\), where \(|A\) is restriction to \(A\), then \(\mathfrak M\)
is a \emph{submodel} of \(\mathfrak M^*\) and \(\mathfrak M^*\) is an
\emph{extension of the first kind} of \(\mathfrak M\).  A bijection
\(\varphi\colon A\to A^*\) that preserves every \(R_k\) in both directions and
commutes with every \(F_k\) makes the models \emph{isomorphic}; \(\mathfrak M\)
is a submodel \emph{in the generalized sense} of \(\mathfrak M^*\) if it is
isomorphic to a submodel of \(\mathfrak M^*\) \cite[pp.~42--43]{Los1955}.
The same notions apply to \(\mathcal E_2\) and \(\mathcal E_3\).  For a model
\(\mathfrak M\) for \(\mathcal E_i\) and \(j<i\), \(\mathfrak M|j\) is the
model for \(\mathcal E_j\) obtained by deleting the interpretations of the
signs not in \(\mathcal E_j\).  A model \(\mathfrak M^*\) for \(\mathcal E_i\)
is an \emph{extension of the second kind} of a model \(\mathfrak M\) for
\(\mathcal E_j\), \(j<i\), if \(\mathfrak M\) is a submodel of
\(\mathfrak M^*|j\), and a \emph{weak extension of the second kind} if
\(\mathfrak M\) is isomorphic to \(\mathfrak M^*|j\)
\cite[p.~43]{Los1955}.
\end{hbdefinition}

Validity of a formula in general passes neither from an extension to a
submodel nor in the other direction \cite[p.~43]{Los1955}; the worked
example of \cref{sec:p03-w03-p03-026-groups} shows both failures for
groups.  Open formulas pass downward.

\begin{hblemma}{thm:p03-w03-p03-026-los55-31}{proved=annex:p03-w03-p03-026-los55-31}{\L o\'s's (3.1): open formulas pass to submodels}
If \(\mathfrak M^*\) is a model of \(\mathcal X\subset\mathcal O_i\) and
\(\mathfrak M\) is a submodel of \(\mathfrak M^*\), then \(\mathfrak M\) is a
model of \(\mathcal X\) \cite[p.~43]{Los1955}.
\end{hblemma}

\subsubsection{Descriptions and the extension theorem (\S\S\,4--5)}

\begin{hbdefinition}{def:p03-w03-p03-026-los55-description}{The description \(\mathcal D(\mathfrak M)\)}
Let \(\mathfrak M=\langle A,R_1,\ldots,F_1,\ldots\rangle\) be a model for
\(\mathcal E_1\), and take the index set of the constants of \(\mathcal E_3\)
to be \(T=A\), so that \(g_a\) names \(a\in A\).  The \emph{description}
\(\mathcal D(\mathfrak M)\) is the set of sentences of \(\mathcal E_3\) of the
forms
(i) \((g_a=g_b)'\) for \(a\neq b\);
(ii) \(r_k(g_{a_1},\ldots,g_{a_{\nu_k}})\) if \(R_k(a_1,\ldots,a_{\nu_k})\);
(iii) \(r_k(g_{a_1},\ldots,g_{a_{\nu_k}})'\) if not
\(R_k(a_1,\ldots,a_{\nu_k})\);
(iv) \(f_k(g_{a_1},\ldots,g_{a_{\varkappa_k}})=g_a\) if
\(F_k(a_1,\ldots,a_{\varkappa_k})=a\)
\cite[pp.~43--44]{Los1955}.  \L o\'s states that the notion is essentially
Robinson's \cite[p.~43]{Los1955}; it is the atomic diagram of
\cref{def:diagrams} written with \L o\'s's constants.
\end{hbdefinition}

\begin{hblemma}{thm:p03-w03-p03-026-los55-desc}{proved=annex:p03-w03-p03-026-los55-desc}{Properties (4.1)--(4.5) of descriptions}
Let \(\mathfrak M\) be a model for \(\mathcal E_1\).
\begin{enumerate}
 \item[(4.1)] \(\mathcal D(\mathfrak M)\subset\mathcal S_3\cap\mathcal O_3\).
 \item[(4.2)] No sign \(q_k\) occurs in a sentence of \(\mathcal D(\mathfrak M)\).
 \item[(4.3)] If \(\mathcal D(\mathfrak M)\subset\mathcal X\subset\mathcal S_3\)
       and \(\mathfrak M^*\) is a model of \(\mathcal X\), then
       \(\mathfrak M^*|1\) is an extension of the first kind and
       \(\mathfrak M^*|2\) an extension of the second kind of
       \(\mathfrak M\), in the generalized sense, through the map
       \(a\mapsto G^*_a\).
 \item[(4.4)] If \(\mathfrak M\) is a model of \(\mathcal X\subset\mathcal E_i\)
       (\(i=1\), or \(i=2\) with \(\mathfrak M\) a model for \(\mathcal E_2\)),
       then \(\mathcal X\cup\mathcal D(\mathfrak M)\) is consistent in
       \(\mathcal E_3\).
 \item[(4.5)] If \(\mathfrak M\) is a model of \(\mathcal A\subset\mathcal E_1\),
       \(\mathcal C\subset\mathcal E_2\), and
       \(\mathcal A\cup\mathcal C\cup\mathcal D(\mathfrak M)\) is
       inconsistent in \(\mathcal E_3\), then some \(\alpha_1\in\mathcal O_1\)
       belongs to \(\mathit{Cn}_2(\mathcal A\cup\mathcal C)\) and is not valid
       in \(\mathfrak M\)
       \cite[p.~44]{Los1955}.
\end{enumerate}
\end{hblemma}

\begin{hbdefinition}{def:p03-w03-p03-026-los55-extension-with-conditions}{Extension with primary and secondary conditions}
Let \(\mathcal A\subset\mathcal E_1\), \(\mathcal C\subset\mathcal E_2\), and
let \(\mathfrak M\) be a model of \(\mathcal A\).  An extension of the second
kind of \(\mathfrak M\) that is a model of \(\mathcal A\cup\mathcal C\) is an
\emph{extension of \(\mathfrak M\) with primary conditions \(\mathcal A\) and
secondary conditions \(\mathcal C\)} \cite[p.~45]{Los1955}.
\end{hbdefinition}

\begin{hbtheorem}{thm:p03-w03-p03-026-los55-thm1}{proved=annex:p03-w03-p03-026-los55-thm1}{\L o\'s's Theorem~1: the extension theorem}
An extension of \(\mathfrak M\) with primary conditions \(\mathcal A\) and
secondary conditions \(\mathcal C\) exists if and only if \(\mathfrak M\) is a
model of \(\mathcal O_1\cap\mathit{Cn}_2(\mathcal A\cup\mathcal C)\)
\cite[p.~45]{Los1955}.
\end{hbtheorem}

\begin{hblemma}{thm:p03-w03-p03-026-los55-remark1}{proved=annex:p03-w03-p03-026-los55-remark1}{Remark~1: secondary conditions without new signs}
For \(\mathcal X\subset\mathcal E_1\),
\[
 \mathcal O_1\cap\mathit{Cn}_1(\mathcal X)=\mathcal O_1\cap\mathit{Cn}_2(\mathcal X).
\]
Consequently, if \(\mathcal C\subset\mathcal E_1\), then \(\mathfrak M\) has an
extension of the first kind that is a model of \(\mathcal A\cup\mathcal C\)
if and only if \(\mathfrak M\) is a model of
\(\mathcal O_1\cap\mathit{Cn}_1(\mathcal A\cup\mathcal C)\)
\cite[p.~45]{Los1955}.
\end{hblemma}

\begin{hbtheorem}{thm:p03-w03-p03-026-los55-thm2}{proved=annex:p03-w03-p03-026-los55-thm2}{\L o\'s's Theorem~2: weak extensions under open secondary conditions}
If \(\mathcal C\subset\mathcal O_2\), then a weak extension of the second kind
of \(\mathfrak M\) that is a model of \(\mathcal A\cup\mathcal C\) exists if
and only if \(\mathfrak M\) is a model of
\(\mathcal O_1\cap\mathit{Cn}_2(\mathcal A\cup\mathcal C)\)
\cite[p.~45]{Los1955}.
\end{hbtheorem}

\L o\'s's Remark~2 explains why Theorem~2 needs \(\mathcal E_2\) to add only
relation signs.  A new \(s\)-ary function sign can be replaced by an
\((s+1)\)-ary relation sign together with formulas saying that the relation
is a function of its last argument.  Those formulas contain quantifiers, so
they belong to \(\mathcal E_2\) but not to \(\mathcal O_2\).  After the
replacement Theorem~1 still applies, but Theorem~2 does not, even when the
original secondary conditions were open \cite[pp.~45--46]{Los1955}.

\subsubsection{\(\mathcal O\)-systems and definable classes (\S\S\,6--8)}

\L o\'s calls \(\mathcal X\subset\mathcal E_1\) a \emph{system} if
\(\mathit{Cn}_1(\mathcal X)=\mathcal X\), and an \emph{\(\mathcal O\)-system}
if moreover \(\mathcal X=\mathit{Cn}_1(\mathcal O_1\cap\mathcal X)\).  He uses
the completeness theorem in the form: if \(\alpha\notin\mathit{Cn}_1(\mathcal X)\),
some model of \(\mathcal X\) does not validate \(\alpha\)
\cite[p.~46]{Los1955}; this is item~1 of
\cref{thm:p03-w03-p03-026-los55-cn}.  By (2.1) and (3.1), every submodel of a
model of an \(\mathcal O\)-system \(\mathcal X\) is a model of \(\mathcal X\);
this is (6.1).  \L o\'s states that the results of his \S\,6 also appear in
his 1955 paper on the algebraic treatment of the methodology of elementary
deductive systems \cite[p.~46]{Los1955}.

\begin{hbtheorem}{thm:p03-w03-p03-026-los55-thm3}{proved=annex:p03-w03-p03-026-los55-thm3}{\L o\'s's Theorem~3: characterization of \(\mathcal O\)-systems}
If every submodel of a model of the system \(\mathcal X\) is a model of
\(\mathcal X\), then \(\mathcal X\) is an \(\mathcal O\)-system
\cite[p.~46]{Los1955}.
\end{hbtheorem}

A class \(\mathbf A\) of models is \emph{elementarily definable} if some
\(\mathcal A\subset\mathcal E_1\) has: \(\mathfrak M\in\mathbf A\) iff
\(\mathfrak M\) is a model of \(\mathcal A\); \(\mathcal A\) is then a
\emph{set of axioms} of \(\mathbf A\) \cite[p.~46]{Los1955}.  \L o\'s
assumes that \(\mathcal E_1\) contains only the signs essential to the
problem at hand, and notes that the choice of signs can matter, as the
groups example of \cref{sec:p03-w03-p03-026-groups} shows
\cite[p.~47]{Los1955}.  For a class \(\mathbf A\), an
\emph{\(\mathbf A\)-submodel} of \(\mathfrak M\) is a submodel of
\(\mathfrak M\) that belongs to \(\mathbf A\).  For
\(\mathcal A\subset\mathcal E_1\), \(\mathbf M(\mathcal A)\) is the class of
all models of \(\mathcal A\); if \(\mathcal A\subset\mathcal B\) then
\(\mathbf M(\mathcal B)\subset\mathbf M(\mathcal A)\), and
\(\mathbf M(\mathcal A)=\mathbf M(\mathcal B)\) iff
\(\mathit{Cn}_1(\mathcal A)=\mathit{Cn}_1(\mathcal B)\), by G\"odel's theorem
\cite[p.~47]{Los1955}.  The last equivalence is immediate from item~1 of
\cref{thm:p03-w03-p03-026-los55-cn}: the consequences of a set are the
formulas valid in all its models.  \L o\'s records that Tarski calls classes
with a finite set of axioms arithmetical classes, a name he avoids because the
notion has little to do with arithmetic \cite[p.~47]{Los1955}.

\begin{hbtheorem}{thm:p03-w03-p03-026-los55-thm45}{proved=annex:p03-w03-p03-026-los55-thm45}{\L o\'s's Theorems~4 and~5}
\begin{enumerate}
 \item[(Theorem~4)] If every submodel of a model belonging to a definable
       class \(\mathbf A\) belongs to \(\mathbf A\), then \(\mathbf A\) has a
       set of axioms contained in \(\mathcal O_1\).
 \item[(Theorem~5)] If \(\mathbf A=\mathbf M(\mathcal A)\) is a definable
       class, \(\mathbf B\) is a definable subclass of \(\mathbf A\), and
       every \(\mathbf A\)-submodel of a model in \(\mathbf B\) belongs to
       \(\mathbf B\), then there is \(\mathcal B\subset\mathcal O_1\) with
       \(\mathbf B=\mathbf M(\mathcal A\cup\mathcal B)\)
       \cite[p.~47]{Los1955}.
\end{enumerate}
\end{hbtheorem}

\L o\'s notes that Theorem~5 generalizes Theorem~3 and can be proved by
repeating that proof \cite[p.~47]{Los1955}; that is the route of the annex.
Theorem~5 is the relative \L o\'s--Tarski theorem
\cref{thm:p03-w03-p03-026-los-tarski-relative}, with open formulas in place
of their universal closures.

\begin{hbcorollary}{thm:p03-w03-p03-026-los55-81}{proved=annex:p03-w03-p03-026-los55-81}{\L o\'s's (8.1)}
If \(\mathbf A=\mathbf M(\mathcal A)\) is a definable class, \(\alpha\) is a
sentence, and every \(\mathbf A\)-submodel of a model in
\(\mathbf M(\mathcal A\cup(\alpha))\) belongs to that class, then
\(\mathit{Cn}_1(\mathcal A)\) contains an equivalence \(\alpha\equiv\beta\)
with \(\beta\) a general sentence
\(\prod_{x_{s_1}}\prod_{x_{s_2}}\cdots\prod_{x_{s_n}}\gamma\),
\(\gamma\in\mathcal O_1\) \cite[p.~48]{Los1955}.
\end{hbcorollary}

\subsubsection{Persistent sentences and submodel classes (\S\S\,9--10)}

A sentence \(\alpha\) is \emph{persistent} with respect to a class
\(\mathbf A\) if, whenever \(\alpha\) is valid in some
\(\mathfrak M\in\mathbf A\), it is valid in every extension of
\(\mathfrak M\) that belongs to \(\mathbf A\) \cite[p.~48]{Los1955}.

\begin{hbtheorem}{thm:p03-w03-p03-026-los55-thm6}{proved=annex:p03-w03-p03-026-los55-thm6}{\L o\'s's Theorem~6: persistent sentences}
If \(\alpha\) is persistent with respect to the definable class
\(\mathbf A=\mathbf M(\mathcal A)\), then \(\mathit{Cn}_1(\mathcal A)\)
contains an equivalence \(\alpha\equiv\beta\) with \(\beta\) an existential
sentence \(\sum_{x_{s_1}}\sum_{x_{s_2}}\cdots\sum_{x_{s_n}}\gamma\),
\(\gamma\in\mathcal O_1\) \cite[p.~48]{Los1955}.
\end{hbtheorem}

The converse of Theorem~6 is item~1 of
\cref{thm:p03-w03-p03-026-los-tarski-sentence} with \(S=\mathcal A\).  \L o\'s
then restates a theorem of Robinson, which he attributes to Robinson's 1951
paper on axiomatic systems with finite models and does not prove: every
sentence persistent with respect to the class of Abelian groups that holds in
some Abelian group holds in some finite group, whether groups are taken with
one or with two primitive operations.  Using definability of Abelian groups,
G\"odel's theorem and Theorem~6, he gives it the two forms (9.1): an
existential sentence consistent with the axioms of Abelian groups is valid in
some finite group; and (9.2): an open formula valid in all finite Abelian
groups is valid in all Abelian groups.  He reports analogous theorems for
commutative rings, also from Robinson, and, on information from
B.~H.~Neumann, that the statement fails for arbitrary groups
\cite[p.~48]{Los1955}.  None of these is proved in the paper, and they are
recorded here only as \L o\'s's report.

For a class \(\mathbf A\), \(\operatorname{sm}(\mathbf A)\) is the class of
all submodels of models in \(\mathbf A\).  Every model is a submodel of
itself, so \(\mathbf A\subset\operatorname{sm}(\mathbf A)\), and
\(\operatorname{sm}(\operatorname{sm}(\mathbf A))=\operatorname{sm}(\mathbf A)\),
because a submodel of a submodel is a submodel; this is (10)
\cite[pp.~48--49]{Los1955}.  In this notation Theorem~4 says that a
definable \(\mathbf A\) satisfies \(\mathbf A=\operatorname{sm}(\mathbf A)\)
iff \(\mathbf A=\mathbf M(\mathcal A)\) for some
\(\mathcal A\subset\mathcal O_1\), and Theorem~5 says that a definable
subclass \(\mathbf B\) of \(\mathbf A=\mathbf M(\mathcal A)\) satisfies
\(\mathbf A\cap\operatorname{sm}(\mathbf B)=\mathbf B\) iff
\(\mathbf B=\mathbf M(\mathcal A\cup\mathcal B)\) for some
\(\mathcal B\subset\mathcal O_1\) \cite[p.~49]{Los1955}.  For a class
\(\mathbf A\) of models for \(\mathcal E_2\), \(\mathbf A|1\) is the class of
the models \(\mathfrak M|1\), \(\mathfrak M\in\mathbf A\).  Because
definable classes are closed under isomorphism, submodels are taken in the
generalized sense whenever \(\operatorname{sm}\) is compared with a definable
class.

\begin{hbtheorem}{thm:p03-w03-p03-026-los55-thm789}{proved=annex:p03-w03-p03-026-los55-thm789}{\L o\'s's Theorems~7, 8 and~9}
\begin{enumerate}
 \item[(Theorem~7)] If \(\mathbf A=\mathbf M(\mathcal A)\) is definable, then
       \(\operatorname{sm}(\mathbf A)\) is definable and
       \(\operatorname{sm}(\mathbf A)=\mathbf M(\mathcal O_1\cap\mathit{Cn}_1(\mathcal A))\).
 \item[(Theorem~8)] If \(\mathbf A=\mathbf M(\mathcal A)\) with
       \(\mathcal A\subset\mathcal E_2\), then
       \(\operatorname{sm}(\mathbf A|1)\) is definable and
       \(\operatorname{sm}(\mathbf A|1)=\mathbf M(\mathcal O_1\cap\mathit{Cn}_2(\mathcal A))\).
 \item[(Theorem~9)] If \(\mathbf A=\mathbf M(\mathcal A)\),
       \(\mathcal A\subset\mathcal E_2\), and
       \(\mathbf A|1\subset\mathbf B=\mathbf M(\mathcal B)\),
       \(\mathcal B\subset\mathcal E_1\), then
       \(\mathbf B\cap\operatorname{sm}(\mathbf A|1)\) is definable and
       \(\mathbf B\cap\operatorname{sm}(\mathbf A|1)
       =\mathbf M[\mathcal B\cup(\mathcal O_1\cap\mathit{Cn}_2(\mathcal A\cup\mathcal B))]\)
       \cite[p.~49]{Los1955}.
\end{enumerate}
\end{hbtheorem}

The display of Theorem~9 in the source reads
\(\mathit{Cn}_2(\mathcal A\cup\mathcal C)\), with a set \(\mathcal C\) that its
hypotheses do not introduce \cite[p.~49]{Los1955}; the statement above reads
it as \(\mathcal B\).  Under the hypothesis \(\mathbf A|1\subset\mathbf B\),
every member of \(\mathcal B\) is valid in all models of \(\mathcal A\), so
\(\mathit{Cn}_2(\mathcal A\cup\mathcal B)=\mathit{Cn}_2(\mathcal A)\) and the
right side is also
\(\mathbf M[\mathcal B\cup(\mathcal O_1\cap\mathit{Cn}_2(\mathcal A))]\).

\subsubsection{Sentences with relation quantifiers (\S\S\,11--12)}

\L o\'s next admits non-elementary sentences of the form
\begin{equation*}
 (*)\qquad \alpha=\sum_{q_1}\sum_{q_2}\cdots\sum_{q_s}\beta ,
\end{equation*}
where \(\beta\in\mathcal S_2\) and no \(q_j\) with \(j>s\) occurs in
\(\beta\).  The signs \(q_j\) in \(\alpha\) range over all relations of the
right arity, while the \(r_k\) and \(f_k\) are fixed.  Such \(\alpha\) is
\emph{valid} in a model
\(\mathfrak M=\langle A,R_1,\ldots,F_1,\ldots\rangle\) of
\(\mathcal A\subset\mathcal E_1\) if there are relations
\(Q_1,\ldots,Q_s\) on \(A\) for which \(\beta\) is valid in
\(\langle A,R_1,\ldots,F_1,\ldots,Q_1,\ldots,Q_s\rangle\); equivalently,
some weak extension of the second kind of \(\mathfrak M\) validates
\(\beta\) \cite[p.~49]{Los1955}.  \(\mathbf M(\mathcal A,\alpha)\) is the class
of models of \(\mathcal A\) in which \(\alpha\) is valid, and for \(\alpha\) of
the form \((*)\)
\begin{equation*}
 (11.1)\qquad \mathbf M(\mathcal A,\alpha)=\mathbf M(\mathcal A\cup(\beta))|1
\end{equation*}
\cite[p.~50]{Los1955}, where on the right the models for \(\mathcal E_2\)
interpret the signs \(q_j\), \(j>s\), arbitrarily.

\begin{hbtheorem}{thm:p03-w03-p03-026-los55-thm1011}{proved=annex:p03-w03-p03-026-los55-thm1011}{\L o\'s's Theorems~10 and~11}
Let \(\alpha=\sum_{q_1}\cdots\sum_{q_s}\beta\) be of the form \((*)\) and
\(\mathcal A\subset\mathcal E_1\).
\begin{enumerate}
 \item[(Theorem~10)] If \(\beta\) is a general sentence
       \(\prod_{x_1}\prod_{x_2}\cdots\prod_{x_n}\gamma\) with
       \(\gamma\in\mathcal O_2\), then \(\mathbf M(\mathcal A,\alpha)\) is
       elementarily definable and
       \(\mathbf M(\mathcal A,\alpha)=\mathbf M\bigl(\mathcal A\cup(\mathcal O_1\cap\mathit{Cn}_2(\mathcal A\cup(\beta)))\bigr)\).
 \item[(Theorem~11)] For arbitrary \(\beta\), the classes
       \(\operatorname{sm}(\mathbf M(\mathcal A,\alpha))\) and
       \(\mathbf M(\mathcal A)\cap\operatorname{sm}(\mathbf M(\mathcal A,\alpha))\)
       are elementarily definable, with
       \(\operatorname{sm}(\mathbf M(\mathcal A,\alpha))=\mathbf M\bigl(\mathcal O_1\cap\mathit{Cn}_2(\mathcal A\cup(\beta))\bigr)\)
       and
       \(\mathbf M(\mathcal A)\cap\operatorname{sm}(\mathbf M(\mathcal A,\alpha))=\mathbf M\bigl(\mathcal A\cup(\mathcal O_1\cap\mathit{Cn}_2(\mathcal A\cup(\beta)))\bigr)\)
       \cite[p.~50]{Los1955}.
\end{enumerate}
\end{hbtheorem}

\begin{hbtheorem}{thm:p03-w03-p03-026-los55-arithmetic}{proved=annex:p03-w03-p03-026-los55-arithmetic}{\L o\'s's \S\,12: a class \(\mathbf M(\mathcal A,\alpha)\) that is not elementarily definable}
Interpret \(f_1,f_2\) as addition and multiplication of natural numbers
(\(\varkappa_1=\varkappa_2=2\)) and \(f_3,f_4\) as the numbers \(0\) and \(1\)
(\(\varkappa_3=\varkappa_4=0\)), and let \(\mathcal A\) be the set of all
formulas valid in this interpretation, a complete system.  Let
\begin{equation*}
 \alpha=\sum_{q_1}\Bigl\{q_1(f_3)\land\prod_{x_1}\bigl[q_1(x_1)\to q_1\bigl(f_1(x_1,f_4)\bigr)\bigr]\land\sum_{x_1}q_1(x_1)'\Bigr\}.
\end{equation*}
Then \(\mathbf M(\mathcal A,\alpha)\) is the class of models of \(\mathcal A\)
in which the (non-elementary) induction axiom fails; it is nonempty, differs
from \(\mathbf M(\mathcal A)\), and is not elementarily definable
\cite[p.~50]{Los1955}.
\end{hbtheorem}

The source states the conclusion and gives its reason, the completeness of
\(\mathcal A\); the annex proof also supplies the two facts left implicit
there, that the class is nonempty and that a model of \(\mathcal A\)
satisfies induction iff it is isomorphic to the natural numbers.

The example bears on Theorem~10: for \(\alpha\) of the form \((*)\) with
\(\beta\) not general, \(\mathbf M(\mathcal A,\alpha)\) need not be
elementarily definable, whereas Theorem~11 shows that its submodel class
always is.

\subsubsection{Algebraic applications (\S\S\,13--16)}

\emph{Semigroup to group} (\S\,13).  Let \(\mathcal E_1\) have one binary
operation \(\cdot\), let \(\mathcal A_1\) consist of the associative law
\((x_1\cdot x_2)\cdot x_3=x_1\cdot(x_2\cdot x_3)\), and let \(\mathcal C_1\)
consist of
\(\prod_{x_1}\prod_{x_2}\sum_{x_3}x_1\cdot x_3=x_2\) and
\(\prod_{x_1}\prod_{x_2}\sum_{x_3}x_3\cdot x_1=x_2\).  Models of
\(\mathcal A_1\) are semigroups and models of
\(\mathcal A_1\cup\mathcal C_1\) are groups \cite[pp.~50--51]{Los1955}; the
second claim is item~1 of \cref{thm:p03-w03-p03-026-groups}, whose two
sentences are equivalent to \(\mathcal A_1\cup\mathcal C_1\).  Extending a
semigroup to a group is therefore the extension problem with primary
conditions \(\mathcal A_1\) and secondary conditions \(\mathcal C_1\).

\begin{hbcorollary}{thm:p03-w03-p03-026-los55-semigroup}{proved=annex:p03-w03-p03-026-los55-semigroup}{Semigroups that extend to groups}
A semigroup \(\mathfrak M\) has an extension of the first kind that is a
group iff \(\mathfrak M\) is a model of
\(\mathcal O_1\cap\mathit{Cn}_1(\mathcal A_1\cup\mathcal C_1)\), a set of open
formulas in the multiplication sign alone.
\end{hbcorollary}

\L o\'s observes that Mal'cev solved the problem in exactly this form, by
open formulas in the multiplication sign, in the first of his two papers
\emph{On the inclusion of associative systems in groups} (Mat.\ Sbornik 6,
1939, and 8, 1940), and proved in the second that any such set of open
formulas must be infinite \cite[p.~51]{Los1955}.  Neither of Mal'cev's
results is proved here.

\emph{Semigroup to ring} (\S\,14).  Let \(\mathcal E_2\) add a binary
operation \(+\), and let \(\mathcal C_2\) consist of
\(\prod_{x_1}\prod_{x_2}\sum_{x_3}x_1+x_3=x_2\), the associative and
commutative laws of \(+\), and the two distributive laws
\(x_1\cdot(x_2+x_3)=x_1\cdot x_2+x_1\cdot x_3\) and
\((x_1+x_2)\cdot x_3=x_1\cdot x_3+x_2\cdot x_3\).  Extending with primary
conditions \(\mathcal A_1\) and secondary conditions \(\mathcal C_2\) is
embedding a semigroup in the multiplicative semigroup of a ring.  By
Theorem~1 it has a solution by quantifier-free conditions; \L o\'s writes
that, as far as he knows, no solution in this or any other form had been
found \cite[p.~51]{Los1955}.  Writing \(\mathbf P\) for the class of models
of \(\mathcal A_1\) and \(\mathcal C_2\) (printed
\(\mathbf M(\mathcal A_1+\mathcal C_2)\)), the problem asks for axioms of the
class of semigroups that are submodels of reducts of members of
\(\mathbf P\), which Theorem~8 gives as
\(\mathcal O_1\cap\mathit{Cn}_2(\mathcal A_1\cup\mathcal C_2)\)
\cite[p.~51]{Los1955}.  In logical terms it asks for a (preferably simple)
\(\mathcal X\subset\mathcal E_1\) with
\(\mathit{Cn}_1(\mathcal X)=\mathit{Cn}_1(\mathcal O_1\cap\mathit{Cn}_2(\mathcal A_1\cup\mathcal C_2))\)
\cite[p.~52]{Los1955}.  The source prints the class as
\(\operatorname{sm}(\mathbf P)|1\), the axiom set with \(\mathcal A\) in place
of \(\mathcal A_1\cup\mathcal C_2\), and the last display with
\(\mathcal A_1\cap\mathcal C_2\); the readings above are the ones under which
Theorem~8 applies, since \(\mathcal A_1\cap\mathcal C_2\) is empty.  \L o\'s
adds that reducing the algebraic problem to a logical one may have no
practical value in general, and points to \S\,16 as a case where the
translation yields the solution \cite[p.~52]{Los1955}.

\emph{Group to ordered group} (\S\,15).  A group is \emph{ordered} by a
relation \(<\) if \(a<b\) implies \(cad<cbd\) for all
\(a,b,c,d\) \cite[p.~52]{Los1955}.  \L o\'s lets \(\gamma\) be the
conjunction of
\(q_1(x_1,x_2)\to q_1(x_2,x_1)'\),
\(q_1(x_3,x_4)\land q_1(x_4,x_5)\to q_1(x_3,x_5)\) and
\(q_1(x_6,x_7)\lor q_1(x_7,x_6)\lor x_6=x_7\), puts
\(\beta=\prod_{x_1}\cdots\prod_{x_7}\gamma\) and
\(\alpha=\sum_{q_1}\beta\), and states that the orderable groups form the
class \(\mathbf U=\mathbf M(\mathcal A_1\cup\mathcal C_1,\alpha)\); since
\(\beta\) is general, Theorem~10 makes \(\mathbf U\) elementarily definable,
with axioms obtained by adding some \(\mathcal X\subset\mathcal O_1\) to
\(\mathcal A_1\cup\mathcal C_1\) \cite[p.~52]{Los1955}.  The three printed
clauses say only that \(q_1\) is a strict linear order.  Every set carries a
strict linear order (by the well-ordering theorem, which uses the Axiom of
Choice), so with \(\gamma\) as
printed \(\alpha\) holds in every group.  The compatibility condition must be
added as a fourth open clause, for instance
\(q_1(x_1,x_2)\to q_1\bigl((x_8\cdot x_1)\cdot x_9,\,(x_8\cdot x_2)\cdot x_9\bigr)\), and
then \(\beta\) remains general; see
\cref{thm:p03-w03-p03-026-los55-ordered}.  \L o\'s adds that the known
theorems on ordered groups produce such a set \(\mathcal X\), citing his 1954
note on the existence of linear order in a group, and that \(\mathcal X\)
must be infinite \cite[p.~52]{Los1955}; these two claims are reported, not
proved, here.

\begin{hbcorollary}{thm:p03-w03-p03-026-los55-ordered}{proved=annex:p03-w03-p03-026-los55-ordered}{Orderable groups form an elementary class}
Let \(\gamma^+\) be the conjunction of the three clauses above and
\(q_1(x_1,x_2)\to q_1\bigl((x_8\cdot x_1)\cdot x_9,\,(x_8\cdot x_2)\cdot x_9\bigr)\), let
\(\beta^+=\prod_{x_1}\cdots\prod_{x_9}\gamma^+\) and
\(\alpha^+=\sum_{q_1}\beta^+\).  A group is orderable iff \(\alpha^+\) is valid
in it, and the class of orderable groups is
\(\mathbf M\bigl(\mathcal A_1\cup\mathcal C_1\cup(\mathcal O_1\cap\mathit{Cn}_2(\mathcal A_1\cup\mathcal C_1\cup(\beta^+)))\bigr)\).
\end{hbcorollary}

\emph{Ternary semigroups} (\S\,16).  Let \(\mathcal E_1\) have one ternary
operation \((a,b,c)\), and let \(\mathcal A_3\) be its complete associativity,
\(((a,b,c),d,e)=(a,(b,c,d),e)=(a,b,(c,d,e))\).  Models of \(\mathcal A_3\) are
\emph{ternary semigroups}, and \(\mathbf T=\mathbf M(\mathcal A_3)\).  A
ternary semigroup is \emph{reducible} if some associative binary operation
satisfies \((a,b,c)=(a\cdot b)\cdot c\) \cite[pp.~52--53]{Los1955}.  \L o\'s
reports, from information of A.~Alexiewicz, that Banach showed during the war
that not every ternary semigroup is reducible, and asked whether every
ternary semigroup extends to a reducible one.  With \(\mathcal E_2\) adding a
binary \(\cdot\) and \(\mathcal C_3\) consisting of
\((a\cdot b)\cdot c=a\cdot(b\cdot c)\) and \((a,b,c)=(a\cdot b)\cdot c\), the
question is whether
\(\operatorname{sm}(\mathbf M(\mathcal A_3\cup\mathcal C_3)|1)=\mathbf T\).
Theorem~8 gives
\(\operatorname{sm}(\mathbf M(\mathcal A_3\cup\mathcal C_3)|1)=\mathbf M(\mathcal O_1\cap\mathit{Cn}_2(\mathcal A_3\cup\mathcal C_3))\).
\L o\'s argues that \(\mathcal A_3\cup\mathcal C_3\) allows only rebracketing,
which \(\mathcal A_3\) already allows for formulas of \(\mathcal O_1\), so that
\(\mathcal O_1\cap\mathit{Cn}_1(\mathcal A_3)=\mathcal O_1\cap\mathit{Cn}_2(\mathcal A_3\cup\mathcal C_3)\),
and hence the submodel class is
\(\mathbf M(\mathcal O_1\cap\mathit{Cn}_1(\mathcal A_3))=\mathbf M(\mathcal A_3)=\mathbf T\):
Banach's problem has a positive solution.  For the middle equality he
sketches only a route through a bracket-free notation for products
\(x_1\cdot x_2\cdots x_n\) \cite[p.~53]{Los1955}.  The annex proves it by a
different route, an explicit reducible extension.

\begin{hbtheorem}{thm:p03-w03-p03-026-los55-ternary}{proved=annex:p03-w03-p03-026-los55-ternary}{Banach's problem: every ternary semigroup extends to a reducible one}
Every ternary semigroup is a submodel of the ternary reduct of a semigroup
\(\langle S,\cdot\rangle\) with \((u,v,w)=(u\cdot v)\cdot w\).  Consequently
\(\mathcal O_1\cap\mathit{Cn}_1(\mathcal A_3)=\mathcal O_1\cap\mathit{Cn}_2(\mathcal A_3\cup\mathcal C_3)\)
and \(\operatorname{sm}(\mathbf M(\mathcal A_3\cup\mathcal C_3)|1)=\mathbf T\).
\end{hbtheorem}

\subsubsection{Weak extensions and \(\omega\)-completeness (\S\,17)}

For weak extensions of the second kind with arbitrary primary and secondary
conditions, \L o\'s leaves the problem open.  The example of \S\,12 suggests
that validity of sentences of \(\mathcal E_1\) in \(\mathfrak M\) cannot be
sufficient, and he proposes that the conditions be sought on the system
\(\mathit{Cn}_3(\mathcal A\cup\mathcal C\cup\mathcal D(\mathfrak M))\) and not
only on systems inside \(\mathcal E_1\) and \(\mathcal E_2\)
\cite[p.~53]{Los1955}.  A system \(\mathcal X\) in \(\mathcal E_3\) is
\emph{\(\omega\)-complete} if, whenever all sentences \(\alpha(g_t)\),
\(t\in T\), belong to \(\mathcal X\), the sentence
\(\prod_{x_s}\alpha(x_s)\) also belongs to \(\mathcal X\), where \(x_s\) does
not occur in \(\alpha(g_t)\) \cite[p.~54]{Los1955}.  \L o\'s states, without
proof, that if an \(\omega\)-complete system containing
\(\mathcal A\cup\mathcal C\cup\mathcal D(\mathfrak M)\) exists in
\(\mathcal E_3\), then \(\mathfrak M\) has a weak extension with conditions
\(\mathcal A\) and \(\mathcal C\), and that the converse holds too
\cite[pp.~53--54]{Los1955}.  He presents this not as a solution but to
indicate where one should be sought.

The forward statement needs the system to be consistent, since the
inconsistent system \(\mathcal E_3\) is \(\omega\)-complete and contains every
set.  The handbook proves the converse and the forward direction in two
cases; the forward direction for a consistent \(\omega\)-complete system that
is neither complete nor countable is proved neither in the source nor here.

\begin{hbtheorem}{thm:p03-w03-p03-026-los55-omega}{proved=annex:p03-w03-p03-026-los55-omega}{\L o\'s's \S\,17 theorem: the proved cases}
Let \(\mathfrak M=\langle A,\ldots\rangle\) be a model for \(\mathcal E_1\),
\(\mathcal A\subset\mathcal E_1\), \(\mathcal C\subset\mathcal E_2\), and let the
constants of \(\mathcal E_3\) be \(g_a\), \(a\in A\).
\begin{enumerate}
 \item If \(\mathfrak M\) has a weak extension of the second kind that is a
       model of \(\mathcal A\cup\mathcal C\), then some consistent, complete,
       \(\omega\)-complete system in \(\mathcal E_3\) contains
       \(\mathcal A\cup\mathcal C\cup\mathcal D(\mathfrak M)\).
 \item If some consistent \(\omega\)-complete system \(\mathcal X\) in
       \(\mathcal E_3\) contains \(\mathcal A\cup\mathcal C\cup\mathcal D(\mathfrak M)\),
       and either \(\mathcal X\) is complete (it contains \(\sigma\) or
       \(\sigma'\) for every \(\sigma\in\mathcal S_3\)) or \(A\) and the set of
       primitive signs of \(\mathcal E_2\) are countable, then \(\mathfrak M\)
       has a weak extension of the second kind that is a model of
       \(\mathcal A\cup\mathcal C\).
\end{enumerate}
\end{hbtheorem}

\L o\'s closes by stating that the theorems and constructions, and above all
the method, are not entirely his own: the method was introduced by Mal'cev
and by Robinson, and Henkin had recently published closely related theorems
\cite[p.~54]{Los1955}.  A note added in proof points to B.~H.~Neumann's
\emph{An embedding theorem for algebraic systems} (Proc.\ London Math.\
Soc.\ 1954) as containing very similar theorems \cite[p.~54]{Los1955}.

\subsection{Relationships}
\label{sec:p03-w03-p03-026-relationships}

The three proofs share one pattern: a semantic closure condition is turned
into an axiomatization by showing that a structure satisfying the
consequences of the right syntactic shape stands in the relevant relation to
a model of the theory, with compactness supplying that model.  The
\L o\'s--Tarski argument needs one step (\cref{thm:p03-w03-p03-026-embed});
the chain argument alternates two steps \(\omega\) times and then uses the
elementary chain theorem to return to the original structure.  Lyndon's
positivity theorem of \cref{sec:p03-n-a-011-02} follows the same pattern,
with surjective homomorphisms in place of inclusions and a polarity
interpolation theorem in place of the diagram argument
\cite{Lyndon1959}.

The chain theorem is named after Chang and after \L o\'s and Suszko; the
proof given here is the alternating-chain argument, and the substructure
theorem is proved both by the embedding criterion and, in
\cref{sec:p03-w03-p03-026-los1955}, by \L o\'s's extension theorem.
